\documentclass[journal]{IEEEtran}

\usepackage{amsmath,amssymb}
\usepackage{graphicx}
\usepackage{cite}
\usepackage[hidelinks]{hyperref}

\begin{document}

\title{Paper Title Goes Here}

\author{
    First~Author,~\IEEEmembership{Member,~IEEE,}
    Second~Author,~\IEEEmembership{Senior~Member,~IEEE}
    \thanks{First Author is with the Department of Computer Science, University Name, City, Country (e-mail: first.author@example.com).}
    \thanks{Second Author is with the Department of Electrical Engineering, University Name, City, Country.}
    \thanks{Manuscript received Month Day, Year; revised Month Day, Year.}
}

\markboth{IEEE Transactions on Example Systems,~Vol.~XX, No.~X, Month~Year}%
{Author \MakeLowercase{\textit{et al.}}: Paper Title Goes Here}

\maketitle

\begin{abstract}
This document is a starting point for an IEEE Transactions/Journal submission. Replace this abstract with a concise summary (150--250 words) of the problem, approach, and key results of the work.
\end{abstract}

\begin{IEEEkeywords}
keyword one, keyword two, keyword three, keyword four
\end{IEEEkeywords}

\section{Introduction}
\IEEEPARstart{T}{his} section introduces the problem, motivates why it matters, summarizes prior approaches and their limitations, and states the paper's contributions as a short itemized list.

\section{Related Work}
Summarize and position the work relative to the most relevant prior literature, organized by theme rather than chronologically.

\section{Methodology}
Describe the proposed method, model, or system. Use subsections to separate distinct components.

\subsection{Problem Formulation}
State assumptions and formal notation here.

\subsection{Proposed Approach}
Describe the core technical contribution here.

\section{Experiments}
Describe the experimental setup: datasets, baselines, evaluation metrics, and implementation details needed for reproducibility.

\subsection{Results}
Present quantitative results, referencing tables and figures as needed.

\section{Conclusion}
Summarize findings, discuss limitations, and outline directions for future work.

\begin{thebibliography}{9}

\bibitem{ref1}
A. Author and B. Author, ``Title of the referenced paper,'' \emph{Journal Name}, vol. 1, no. 1, pp. 1--10, Year.

\bibitem{ref2}
C. Author, D. Author, and E. Author, ``Title of another referenced paper,'' in \emph{Proc. Conference Name}, Year, pp. 1--6.

\end{thebibliography}

\end{document}
